\documentclass[11pt,a4paper]{article}

\usepackage[margin=2.4cm]{geometry}
\usepackage{amsmath,amssymb}
\usepackage{enumitem}
\usepackage{hyperref}

\title{Mathematical and Algorithmic Description of ND Posture Guard}
\author{}
\date{}

\begin{document}
\maketitle

\section{Purpose}

ND Posture Guard is a personalized, example-based posture classifier. It does
not assume that one fixed geometric posture is correct for every person. The
user supplies examples of two classes,
\[
y \in \{0,1\}, \qquad 0=\mathrm{GOOD},\;1=\mathrm{BAD},
\]
and the program learns how the user's own head--neck--shoulder geometry differs
between those classes.

\section{Seven-landmark representation}

One posture sample is described by seven image points
\[
P=(p_1,\ldots,p_7), \qquad p_i=(x_i,y_i)\in\mathbb{R}^2,
\]
with the semantic order
\begin{align*}
p_1 &: \text{eye center on the left side of the displayed image},\\
p_2 &: \text{eye center on the right side of the displayed image},\\
p_3 &: \text{lowest visible chin point},\\
p_4 &: \text{left-image neck/shoulder junction},\\
p_5 &: \text{left-image outer shoulder endpoint},\\
p_6 &: \text{right-image neck/shoulder junction},\\
p_7 &: \text{right-image outer shoulder endpoint}.
\end{align*}

The first three points define the face triangle. Points $(p_4,p_5)$ define the
left shoulder segment and $(p_6,p_7)$ the right shoulder segment. During
monitoring, these are the same three shapes drawn on the live image, so the
visual overlay shows the geometry actually used by the classifier.

\section{Geometric plausibility}

Training and tracked geometries are rejected when their screen-space ordering is
physically implausible. Examples include an eye order reversal after
canonicalization, a chin above the eyes, shoulder endpoints pointing inward, an
unrealistically short shoulder span, or points too far outside the image. These
checks are not the classifier itself; they are data-quality constraints that
keep malformed annotations from entering the model.

\section{Translation and scale normalization}

Let the two inner shoulder points be $p_4$ and $p_6$. The reference center is
\[
c_s = \frac{p_4+p_6}{2}.
\]
Let the outer-shoulder span be
\[
s = \lVert p_7-p_5\rVert_2.
\]
The implementation lower-bounds $s$ by both $0.10$ of the image width and by
$1$ pixel to avoid division by zero or excessive amplification.

Each raw point becomes
\[
\tilde p_i = \frac{p_i-c_s}{s}, \qquad i=1,\ldots,7.
\]
This gives $14$ normalized coordinate values. Translation of the whole body in
the image is largely removed by subtracting $c_s$, and global image scale is
largely removed by dividing by the shoulder span $s$.

\section{The 28-dimensional posture feature}

Define the eye midpoint
\[
e_m=\frac{p_1+p_2}{2}.
\]
The implementation additionally uses four two-dimensional normalized vectors:
\begin{align*}
v_e &= \frac{p_2-p_1}{s}, &
 v_c &= \frac{p_3-e_m}{s},\\
v_L &= \frac{p_5-p_4}{s}, &
 v_R &= \frac{p_7-p_6}{s}.
\end{align*}
These contribute $8$ more numbers.

Six scalar descriptors are then added:
\begin{align*}
w_e &= \frac{\lVert p_2-p_1\rVert_2}{s},\\
h_f &= \frac{\lVert p_3-e_m\rVert_2}{s},\\
t_o &= \frac{y_7-y_5}{s},\\
t_i &= \frac{y_6-y_4}{s},\\
d_e &= \frac{c_{s,y}-e_{m,y}}{s},\\
d_c &= \frac{c_{s,y}-y_3}{s}.
\end{align*}
Thus
\[
x\in\mathbb{R}^{28}
\]
contains $14+8+6=28$ geometry values.

\section{Robust feature scaling}

Different feature dimensions can have different natural magnitudes. Before
distance calculations, each dimension is robustly standardized. For training
feature matrix $X$, the per-dimension center is
\[
m_j=\operatorname{median}_i X_{ij}.
\]
The median absolute deviation is
\[
\operatorname{MAD}_j=
\operatorname{median}_i\left|X_{ij}-m_j\right|.
\]
The robust scale is
\[
a_j=\max\left(1.4826\,\operatorname{MAD}_j,\;0.012\right).
\]
A feature is transformed as
\[
z_j=\frac{x_j-m_j}{a_j}.
\]
The $0.012$ floor prevents nearly constant dimensions from exploding because of
very small numerical differences.

\section{Nearest-example class distance}

For the scaled runtime vector $z$, Euclidean distances to all scaled training
examples $z_i$ are
\[
r_i=\lVert z-z_i\rVert_2.
\]
For each class, the classifier averages only the nearest few examples. If $k$
is the configured neighbor count, the implementation uses at most three:
\[
k_c=\min(k,N_c,3).
\]
Let $D_G$ be the mean of the $k_G$ nearest GOOD distances and $D_B$ the mean of
the $k_B$ nearest BAD distances.

The distance-based BAD score is
\[
s_D=
\frac{D_G^2}{D_G^2+D_B^2}.
\]
Therefore $s_D\rightarrow 0$ when the current vector is much closer to GOOD and
$s_D\rightarrow 1$ when it is much closer to BAD. If both distances are
numerically zero, the score is defined as $0.5$.

\section{Supervised GOOD-to-BAD posture axis}

The classifier also learns an explicit direction from the GOOD cloud toward the
BAD cloud. Let $G$ and $B$ be the scaled GOOD and BAD samples. Their robust
centers are coordinate-wise medians
\[
\mu_G=\operatorname{median}(G),\qquad
\mu_B=\operatorname{median}(B).
\]
The raw separation is
\[
\Delta=\mu_B-\mu_G.
\]

Dimensions that fluctuate strongly inside the classes should matter less. Let
\[
v=\frac{1}{2}\left(\operatorname{Var}(G)+\operatorname{Var}(B)\right)
\]
be the pooled per-dimension variance. A variance floor $v_0$ is estimated from
the positive pooled variances and is never allowed below $0.08$. The
variance-normalized direction is
\[
d' = \frac{\Delta}{\sqrt{v+v_0}},
\qquad
 d=\frac{d'}{\lVert d'\rVert_2}.
\]
The sign is chosen so that BAD examples project above GOOD examples.

For scalar projections
\[
q_i=z_i^\top d,
\]
the class medians are
\[
q_G=\operatorname{median}_{i\in G}q_i,
\qquad
q_B=\operatorname{median}_{i\in B}q_i,
\]
and the midpoint decision boundary is
\[
b=\frac{q_G+q_B}{2}.
\]
The projection scale used by the code is
\[
\tau=\max\left(0.24(q_B-q_G),\;1.5\eta,\;0.12\right),
\]
where $\eta$ is the 75th percentile of the absolute residual projections around
the two class medians.

For a runtime vector $z$, its projection score is obtained with a logistic
function
\[
s_P=\sigma\left(\frac{z^\top d-b}{\tau}\right),
\qquad
\sigma(u)=\frac{1}{1+e^{-u}}.
\]
The implementation clips the logistic argument to $[-12,12]$ for numerical
stability.

\section{Combining the two classifiers}

When the supervised axis is available, the raw BAD score is
\[
s_{\mathrm{raw}}=0.40s_D+0.60s_P.
\]
If the axis cannot be estimated reliably, the classifier falls back to
\[
s_{\mathrm{raw}}=s_D.
\]
This combination prevents the learned direction from completely overriding the
actual nearest labeled examples while still giving the learned GOOD-to-BAD axis
more weight.

\section{Temporal smoothing and confidence}

The last five raw scores are stored and the runtime BAD score is their median:
\[
s_t=\operatorname{median}
\left(s_{\mathrm{raw},t-4},\ldots,s_{\mathrm{raw},t}\right).
\]
This reduces one-frame noise.

The displayed posture class uses the natural midpoint
\[
s_t>0.5 \Rightarrow \mathrm{BAD},
\qquad
s_t\le 0.5 \Rightarrow \mathrm{GOOD}.
\]
Classification confidence is
\[
C=\operatorname{clip}\left(2|s_t-0.5|,0,1\right).
\]
Thus a score near $0.5$ has low confidence, whereas a score near $0$ or $1$ has
high confidence.

\section{Alert threshold and hysteresis}

The displayed GOOD/BAD boundary and the sound threshold are intentionally
separate. Let $T$ be the user-selected beep threshold. A frame counts as an
alerting frame only when
\[
s_t>T.
\]
The alarm is emitted only after this remains true for the configured number of
consecutive frames. This prevents a brief threshold crossing from generating an
immediate beep.

\section{Tracking in live monitoring}

Once a GOOD/BAD model is ready, the seven image landmarks are initialized from
training templates and followed through the live stream using pyramidal Lucas--
Kanade optical flow. Local template matching is used to recover failed points
and periodically re-anchor them to reduce drift.

The monitoring window draws the exact current estimate used by classification:
\begin{itemize}[leftmargin=2em]
    \item the eye--eye--chin triangle,
    \item the left-image shoulder segment,
    \item the right-image shoulder segment,
    \item the seven corresponding points.
\end{itemize}
The overlay color reflects the current state and tracking reliability.

\section{Why camera-frame UI throttling matters}

The camera worker and the Qt GUI run on different threads. Each displayed camera
frame requires conversion from BGR to RGB and a copied \texttt{QImage}. If the
worker emits frames faster than the GUI can consume them, queued Qt events can
accumulate faster than they are processed. Training has separate immediate
signals for selected points, progress, status, and save completion, so camera
images themselves do not need to be delivered at the full capture rate.

The responsive worker therefore bounds camera-image delivery to at most 8 FPS
both during normal monitoring and while training. This keeps the GUI event queue
bounded without delaying point-selection or training-progress feedback.

\section{Interpretation and limitations}

The model is best interpreted as a personalized 2D geometry similarity model. It
captures features such as head displacement relative to the shoulders, chin
drop, shoulder asymmetry, and changes in head/neck/shoulder configuration. It is
not a 3D biomechanical reconstruction and is not a medical diagnostic system.
Its quality depends on representative GOOD/BAD samples, stable camera placement,
lighting, and reliable landmark tracking.

\end{document}
